\section{A fixed coefficient norm for the diagonal canonical map}
\label{supp:canonical-fixed-norm}

The diagonal response of a canonical slot involves only the linear
conditional mean.  The following estimate separates this fact from the
nonlinear construction.  In particular, the same exponential norm is used
on the two sides of every block transformation.

We use the distance $|r|_1=|r_1|+|r_2|$ on the square lattice.  Fix a
shortest lattice path from $o$ to each $o+r_i$, and count its edges with
multiplicity, including multiplicities shared by different paths.  Our
path parameter and weight are
\begin{equation}
 u(r_1,\ldots,r_n)=1+\sum_{i=1}^n|r_i|_1,
 \qquad W_\sigma(u)=e^{\sigma u}(1+u)^2.
 \label{supp:canonical-path}
\end{equation}
The same convention is imposed at every scale.  The union of these paths
is connected and joins the anchor to every argument.  Thus this is a
specific choice of the comparable joining paths used to define the
canonical coefficients; it does not change their polynomials or their
polynomial supports.  We do not identify norms with different exponential
exponents.

For $3\le n\le6$, fix once and for all a basis of the invariant
color tensors of degree $n$.  For a tensor $T$ in that basis, a label represents
$T(y_o(o+r_1),\ldots,y_o(o+r_n))$, with a scalar coefficient $a$.
Its contribution to the coefficient norm is $|a|W_\sigma(u)$.
Labels with a zero displacement may be discarded, since $y_o(o)=0$.
Separate labels may be
retained even when they represent the same polynomial; combining their
coefficients can only decrease the norm.

Here is the precise quadratic convention.  Let $\mathcal D$ be the
eight symmetries of the square, acting on displacement vectors, and set
\begin{align}
 M_{r,s}(y_o)&=\frac18\sum_{g\in\mathcal D}
          y_o(o+gr)\mathbin{\cdot}y_o(o+gs),\nonumber\\
 S_o(y_o)&=\frac14\sum_{|e|_1=1}|y_o(o+e)|^2,
 \qquad Q_{r,s}=M_{r,s}-(r\mathbin{\cdot}s)S_o.
 \label{supp:quadratic-packet}
\end{align}
On every affine field $y_o(o+r)=Ar$, where $A:\mathbb R^2\to\mathbb R^3$
is linear, $Q_{r,s}$ is zero.  Indeed,
\begin{equation}
 \frac18\sum_{g\in\mathcal D}(gr)_i(gs)_j
       =\frac{r\mathbin{\cdot}s}{2}\delta_{ij},
 \qquad S_o(Ar)=\frac12\sum_{i=1}^2|Ae_i|^2.
 \label{supp:quadratic-affine-cancellation}
\end{equation}
A compensated label is the entire packet $aQ_{r,s}$, including its
tagged compensating coefficients.  We assign each of those coefficients
the weight of its own path.  The packet therefore has norm
\begin{equation}
 |a|\{W_\sigma(u(r,s))
           +|r\mathbin{\cdot}s|W_\sigma(3)\}.
 \label{supp:quadratic-packet-norm}
\end{equation}
In particular, a compensating nearest-neighbor term has path parameter
$3$, not the path parameter of the possibly long original label.
Keeping its tag records the cancellation before absolute values are
taken.  For every degree let $\|A_n\|_\sigma$ be the supremum over anchors
of the sum of the corresponding label norms.

Use block-centered fine coordinates, so the $L^2$ fine anchors $o$ assigned
to a coarse anchor $O$ satisfy $|o/L-O|_1\le1$.  Fine coordinate differences
are integral; their absolute origins may be shifted by a half-integer.
Write $D_i f(x)=f(x+e_i)-f(x)$.  The kernel assumptions needed for the
estimate are the following, with $c_*>0$, $K_1$, and $K_2$ independent of
the history and of all $L\ge L_0$:
\begin{equation}
 \sup_x\sum_z e^{c_*|z-x/L|_1}
       |D_{i_1}\cdots D_{i_k}P(x,z)|
       \le K_kL^{-k},\qquad k=1,2.
 \label{supp:fixed-norm-kernel-assumption}
\end{equation}
For the statement about stiffness extraction we also use the exact
constant and affine identities
\begin{equation}
 \sum_zP(x,z)=1,\qquad \sum_zP(x,z)z=x/L.
 \label{supp:fixed-norm-affine-assumption}
\end{equation}
These identities are separate from the difference estimates.

The linear substitution in a label at $o$ is
\begin{equation}
 y_o(o+r)\longmapsto\sum_z
       \bigl(P(o+r,z)-P(o,z)\bigr)y_O(z),\qquad O=[o].
 \label{supp:fixed-norm-substitution}
\end{equation}
It is the linear part of changing anchor after applying the conditional
mean; the constant identity makes it independent of the chosen constant
coarse field.  Output labels are assigned to $O$ and contributions of all
$o$ in that block are added.  Let $B_O$ denote that raw output polynomial.
For the norm estimate at degree two, define $A'_2$ to be the coefficient
list obtained by averaging $B_O$ under the eight square symmetries about
$O$ and then compensating every orbit as in
\eqref{supp:quadratic-packet}.  These operations preserve the polynomial
under the additional symmetry and affine assumptions below.  At higher
degrees, $A'_n$ is simply the raw output list.

\begin{lemma}[One norm and constants chosen before the block size]
\label{supp:one-norm-contraction}
Fix $0<\sigma\le c_*/2$.  Under
\eqref{supp:fixed-norm-kernel-assumption}, substitution
\eqref{supp:fixed-norm-substitution} has the following bounds in the
single norm just defined:
\begin{equation}
 \|A'_n\|_\sigma\le C L^{2-n}\|A_n\|_\sigma
       \quad(3\le n\le6),\qquad
 \|A'_2\|_\sigma\le \frac C L\|A_2\|_\sigma.
 \label{supp:one-norm-conclusion}
\end{equation}
The quadratic assertion is for sums of compensated packets and includes
the cost of output compensation.  To regard this compensation as a
representation of the same polynomial, assume additionally
\eqref{supp:fixed-norm-affine-assumption} and that the output polynomial
at each coarse anchor is invariant under the square symmetries.  The
latter holds for bulk translation-invariant input coefficients and a
block-centered, square-equivariant kernel $P$.  In this case the entire
transferred quadratic polynomial has zero affine Hessian, and its
stiffness extraction is exactly zero.

The constant $C$ and a lower threshold for $L$ depend only on $\sigma$,
$c_*$, $K_1$, $K_2$, and the fixed tensor conventions.  They do not depend
on $L$, the history, or the coefficient lists.  One can consequently
choose $L$ after $C$ so that every bound in
\eqref{supp:one-norm-conclusion} is a strict contraction.
\end{lemma}

\begin{proof}
Put $\eta=2\sigma$ and, for a scalar row $F$, write
\[
 |F|_{\eta,O}=\sum_z e^{\eta|z-O|_1}|F(z)|.
\]
For $m=|r|_1$, telescoping along a shortest path and using
\eqref{supp:fixed-norm-kernel-assumption} gives
\begin{equation}
 |P(o+r,\cdot)-P(o,\cdot)|_{\eta,O}
       \le C\frac mL e^{\eta m/L}.
 \label{supp:row-first}
\end{equation}
Here and in the next two displays $C$ is independent of $L$.  To check
the exponential factor, every fine point in that telescoping sum lies
within distance $m+1$ of $o$.  For such a point $x$,
$|z-O|_1\le |z-x/L|_1+1+(m+1)/L$; the extra bounded factor is absorbed
into $C$.

There is also the exact discrete first-order decomposition
\begin{equation}
 P(o+r,\cdot)-P(o,\cdot)=\sum_i r_iG_i+E_r,
 \quad G_i=D_iP(o,\cdot),\quad
 |G_i|_{\eta,O}\le C/L,
 \quad |E_r|_{\eta,O}\le C\frac{m^2}{L^2}e^{\eta m/L}.
 \label{supp:row-second}
\end{equation}
For completeness, a positive path edge contributes $D_iP(x,\cdot)$,
and a negative edge contributes $-D_iP(x-e_i,\cdot)$.  Subtract the
corresponding signed $D_iP(o,\cdot)$ at each edge.  The signed sum of
these reference rows is $\sum_i r_iG_i$.  Each remaining difference of
first differences telescopes into at most $m+1$ second differences;
there are $m$ edges.  The bound for $k=2$ proves the last estimate after
enlarging $C$.  The case $r=0$ is identically zero.

We will pass from row estimates to coefficient estimates with the
elementary inequality
\begin{equation}
 W_\sigma(1+t)=e^{\sigma(1+t)}(2+t)^2
       \le C_\sigma e^{2\sigma t},\qquad t\ge0,
 \label{supp:row-to-path}
\end{equation}
where $C_\sigma=e^\sigma\sup_{t\ge0}(2+t)^2e^{-\sigma t}<\infty$.
Thus a tensor product of $n$ rows has output coefficient norm at most
$C_\sigma$ times the product of their $|\cdot|_{\eta,O}$ norms.
No input or output coefficient norm has changed: $\eta$ is solely the
weight used to estimate the kernels.

Consider first a degree-$n$ label, $n\ge3$, and put
$u=1+\sum_i|r_i|_1$.  Equations \eqref{supp:row-first} and
\eqref{supp:row-to-path} bound its output coefficient norm, before the
anchor sum, by
\begin{equation}
 C_\sigma |a| L^{-n}(1+u)^n e^{2\sigma u/L}.
 \label{supp:degree-n-before-anchor}
\end{equation}
The substitution acts on position indices only; it leaves the fixed
color tensor unchanged.  Any chosen canonical reordering of its at
most six arguments costs only a fixed finite-dimensional constant.
For $L\ge4$,
\begin{equation}
 \frac{(1+u)^n e^{2\sigma u/L}}{W_\sigma(u)}
       \le (1+u)^{n-2}e^{-\sigma u/2}\le C_{n,\sigma}.
 \label{supp:spare-width-fixed-space}
\end{equation}
This is the precise spare exponential factor: it comes from rescaling
the old displacement by $L$, and is available on every iteration of
the same space.  Summing the old labels and the $L^2$ anchors gives the
first estimate of \eqref{supp:one-norm-conclusion}.

For a quadratic packet, insert \eqref{supp:row-second} into
\eqref{supp:quadratic-packet}.  The terms containing two $G$ rows cancel
as coefficient arrays, by \eqref{supp:quadratic-affine-cancellation}.
This cancellation precedes taking absolute values.  Every remaining
term contains an $E$ row.  With $u=1+|r|_1+|s|_1$, the two terms with
one $E$ have bounds
\[
 C L^{-3}(|r|_1|s|_1^2+|r|_1^2|s|_1)e^{2\sigma u/L},
\]
and the term with two $E$ rows is at most
$C L^{-4}|r|_1^2|s|_1^2e^{2\sigma u/L}$.
The compensating nearest-neighbor terms obey the same estimate with
an additional factor $|r\cdot s|\le |r|_1|s|_1$ and bounded displacements.
After \eqref{supp:row-to-path}, the entire raw output is therefore at most
\begin{equation}
 C_\sigma |a|L^{-3}(1+u)^4e^{2\sigma u/L}
       \le C_\sigma' L^{-3}|a|W_\sigma(u).
 \label{supp:quadratic-before-anchor}
\end{equation}
The final inequality is \eqref{supp:spare-width-fixed-space} with
$n=4$.  It already uses only the first summand of the input packet norm.

We next verify the cost and meaning of output compensation.  Symmetry
averaging does not increase the coefficient norm, since every element
of an orbit has the same $u$.  A raw orbit with displacements $p,q$
and coefficient $b$ gains the compensator $-b(p\cdot q)S_O$.  Its
additional norm is at most
\begin{equation}
 |b|\,|p\cdot q|W_\sigma(3)
 \le \tfrac14 W_\sigma(3)|b|(u(p,q)-1)^2
 \le \tfrac14 W_\sigma(3)|b|W_\sigma(u(p,q)).
 \label{supp:output-comp-cost}
\end{equation}
Thus compensation costs at most the fixed factor
$1+W_\sigma(3)/4$.  Together with
\eqref{supp:quadratic-before-anchor} and the $L^2$ anchor count this
proves the quadratic norm bound.

Finally suppose the exact affine and symmetry assumptions hold.  On
the coarse affine field $y_O(z)=A(z-O)$,
\eqref{supp:fixed-norm-affine-assumption} gives
\[
 \sum_z\bigl(P(o+r,z)-P(o,z)\bigr)A(z-O)=Ar/L.
\]
Consequently every old packet transfers to a polynomial zero on every
affine field, before the sum over anchors.  Polarization gives zero
affine bilinear Hessian as well.  Express a square-invariant output as
the sum of its raw orbit labels $bM_{p,q}$.  Its affine evaluation is
$\frac12\sum b(p\cdot q)\|A\|_{\mathrm{HS}}^2$, so
$\sum b(p\cdot q)=0$.  Absolute convergence follows from the norm
estimates.  The aggregate compensator is therefore exactly zero;
compensation changes only the labelled representation and the extracted
scalar stiffness is zero.  This proves every assertion.
\end{proof}

\begin{remark}[Symmetry, folding, and the scope of this estimate]
The scalar compensation by $S_O$ is appropriate to bulk square-invariant
densities.  It is not a claim that an arbitrary anchor-dependent array
becomes square invariant under blocking.  The latter would require a
larger space with tensor-valued affine compensation.  In the bulk case,
square covariance of $P$, invariance of the input, and the permutation
of all $L^2$ anchor phases under a block-centered square symmetry imply
the required output invariance.  One must sum those phases before
grouping output orbits.

The short path weight of a compensating coefficient is only a convention
for this unmasked coefficient norm.  Its tag retains the original packet
and cancellation provenance.  It does not authorize shortening any mask,
complete read support, or dependency graph in the nonlinear construction.

All estimates and affine normalizations above concern unfolded
coefficients.  Folding their absolutely convergent lists onto a torus
does not increase the norm if lifted path labels are retained, and also
does not increase it if shorter physical paths are subsequently used.
Affine test fields need not be periodic, since they are used before
folding.  The lemma proves only the diagonal coefficient estimate and
its zero-stiffness response.  Off-diagonal connected sums and the exact
nonlinear comparison require their own estimates.
\end{remark}
